\documentclass[11pt]{article}
\usepackage[margin=1in]{geometry}
\usepackage{amsmath, amssymb}
\usepackage{listings}
\usepackage{xcolor}
\usepackage{hyperref}

\lstset{
  language=Matlab,
  basicstyle=\ttfamily\small,
  keywordstyle=\color{blue},
  commentstyle=\color{gray},
  stringstyle=\color{purple},
  showstringspaces=false,
  frame=single,
  breaklines=true,
  columns=fullflexible,
}

\title{MAT 3510 --- Homework 2 Study Notes\\
\large Gaussian Elimination and LU Decomposition}
\author{}
\date{}

\begin{document}
\maketitle

These notes cover the \emph{math} behind the problem, some guiding
questions, and MATLAB syntax reference. Per the course AI policy, they
deliberately do \textbf{not} lay out the algorithm step by step in code
--- turning the math into loops is the part you do yourself.

\section*{1. The big picture}

Solving $Ax=b$ by Gaussian elimination secretly does two things at once:
it turns $A$ into an upper-triangular matrix $U$, and it records every
row operation it used. Recording those operations in a lower-triangular
matrix $L$ gives the factorization
\[
  A = LU .
\]
Once you have it, $Ax=b$ becomes $L(Ux)=b$, which splits into two
\emph{easy} triangular problems:
\[
  \underbrace{Lz = b}_{\text{forward substitution}}
  \qquad\text{then}\qquad
  \underbrace{Ux = z}_{\text{backward substitution}} .
\]

\textbf{Why bother factoring instead of just eliminating?} Factoring costs
about $\tfrac{2}{3}d^3$ flops; each triangular solve costs only about
$d^2$. If you must solve with the same $A$ and many different $b$'s, you
pay the expensive part once.

\section*{2. What $L$ and $U$ look like}

\begin{itemize}
  \item $U$ is upper triangular: it is exactly the matrix you are left
        with after elimination zeros out everything below the diagonal.
  \item $L$ is \emph{unit} lower triangular: ones on the diagonal, zeros
        above it.
  \item The entries of $L$ below the diagonal are the \textbf{multipliers}.
        If, to zero out entry $(i,k)$, you did
        \[
          \text{row}_i \;\leftarrow\; \text{row}_i - \ell_{ik}\,\text{row}_k,
          \qquad
          \ell_{ik} = \frac{a_{ik}}{a_{kk}},
        \]
        then that number $\ell_{ik}$ is stored as $L(i,k)$.
        (Here $a_{ik}$ and $a_{kk}$ mean the entries \emph{at that moment}
        of the elimination, not the original entries of $A$.)
\end{itemize}

\subsection*{A tiny example by hand ($2\times 2$)}
\[
  A=\begin{pmatrix} 2 & 1\\ 6 & 5\end{pmatrix}.
\]
To kill the $6$, subtract $3\times$ row 1 from row 2 (multiplier
$\ell_{21}=6/2=3$):
\[
  U=\begin{pmatrix} 2 & 1\\ 0 & 2\end{pmatrix},\qquad
  L=\begin{pmatrix} 1 & 0\\ 3 & 1\end{pmatrix},\qquad
  LU=\begin{pmatrix} 2 & 1\\ 6 & 5\end{pmatrix}=A.\ \checkmark
\]
Try a $3\times 3$ by hand before coding. Notice which multipliers you
compute in which order --- that order \emph{is} your loop structure.

\section*{3. Triangular solves}

For an upper-triangular system $Ux=z$, look at the \textbf{last}
equation first: it involves only one unknown,
\[
  u_{dd}\,x_d = z_d .
\]
Once $x_d$ is known, the second-to-last equation has only one
\emph{unknown} left. Write out the general $i$-th equation of $Ux=z$ and
ask: which $x_j$'s are already known when I reach row $i$, and how do I
isolate $x_i$? Forward substitution for $Lz=b$ is the mirror image,
starting from the \emph{first} row. (Bonus: since $L$ has ones on its
diagonal, what division do you get to skip?)

\section*{4. Guiding questions for designing your function}

\begin{enumerate}
  \item \textbf{Inputs/outputs.} What should your function take in and
        return? Just $x$? Or $L$ and $U$ too (useful for part (2) and for
        checking)? Consider separate functions for the factorization and
        for each triangular solve.
  \item \textbf{Columns.} Elimination works column by column. For a
        $d\times d$ matrix, which columns need entries eliminated below
        the diagonal? (Does the last column need any work?)
  \item \textbf{Rows.} For a fixed column $k$, which rows $i$ get a
        multiplier? What is the range of $i$?
  \item \textbf{Which entries change?} When you update row $i$ using
        row $k$, which columns of that row actually change? Can you do
        the update with a single vectorized row operation instead of an
        inner loop?
  \item \textbf{Zero pivots.} What happens to $\ell_{ik}=a_{ik}/a_{kk}$
        if $a_{kk}=0$? Your function should work for a ``general''
        $d$, so at minimum decide what it should do in that case (detect
        it and \texttt{error}, or swap rows --- partial pivoting, which
        gives $PA=LU$). For the given test system, work out the pivots by
        hand: do you ever hit a zero?
  \item \textbf{Counting.} Your loops should cost $O(d^3)$ for the
        factorization and $O(d^2)$ per triangular solve. Does your loop
        nesting depth match that?
\end{enumerate}

\section*{5. How to check your answer}

\begin{itemize}
  \item \textbf{Reconstruction:} \texttt{norm(L*U - A)} should be at
        round-off level (around $10^{-15}$).
  \item \textbf{Structure:} is your $L$ really unit lower triangular and
        $U$ really upper triangular? (\texttt{istril}, \texttt{istriu}
        test this.)
  \item \textbf{Residual:} \texttt{norm(A*x - b)} should be tiny.
  \item \textbf{Against MATLAB:} compare your $x$ to \verb|A\b|.
        \emph{Careful:} MATLAB's built-in \texttt{[L,U,P] = lu(A)} uses
        partial pivoting, so its $L$ and $U$ may differ from yours even
        when both are correct. Compare solutions and reconstructions, not
        factor-by-factor.
\end{itemize}

\section*{6. MATLAB syntax reference}

These are general language features, not the algorithm.

\begin{lstlisting}
% A function file: must be saved as myfunc.m (file name = function name)
function [out1, out2] = myfunc(in1, in2)
    ...
end

d = size(A, 1);        % number of rows of A
I = eye(d);            % d-by-d identity matrix
z = zeros(d, 1);       % d-by-1 column vector of zeros

for k = 1:d-1          % loop k = 1, 2, ..., d-1
    ...
end
for i = d:-1:1         % loop counting DOWN: d, d-1, ..., 1
    ...
end

A(i, :)                % entire row i
A(i, k:end)            % row i, columns k through the last
A(k+1:end, k)          % column k, rows below the diagonal
v(j+1:end)' * w(j+1:end)   % dot product of two sub-vectors

error('Zero pivot at step %d', k);   % stop with a message
A = [1 -2 3; 0 1 -1; 1 3 0];         % matrix literal: ; separates rows
b = [4; -3; 1];                      % COLUMN vector (note the semicolons)
format long                           % show more digits when printing
\end{lstlisting}

\section*{7. Common pitfalls}

\begin{itemize}
  \item \textbf{Row vs.\ column $b$.} \texttt{[4 -3 1 -3]} is a row
        vector; \texttt{[4; -3; 1; -3]} is a column. Mixing them causes
        dimension errors in \texttt{A*x - b}.
  \item \textbf{Using stale entries.} The multiplier uses the
        \emph{current} value of $a_{ik}$, after earlier eliminations have
        modified it --- not the original $A$.
  \item \textbf{Off-by-one ranges.} Most bugs in this problem are loop
        bounds. Trace your code by hand on the $2\times 2$ example above
        and on a $3\times 3$ you've done on paper.
  \item \textbf{File names.} Each function you call from your test
        script must live in its own \texttt{.m} file with the matching
        name (same lesson as HW1).
\end{itemize}

\end{document}
